\chapter{Execution Stage}

\section{Overview}

The Execution (EXE) stage is responsible for performing the actual computation requested by
each instruction. After operands have been read from the physical register file, instructions are
dispatched to the appropriate execution unit based on their operation type.

Unlike a simple in-order processor that contains a single ALU, Sargantana contains multiple
specialized execution units capable of executing integer, branch, multiplication, division,
floating-point, SIMD and memory operations simultaneously.

The entire execution subsystem is implemented under:

\begin{center}
\rtldir{rtl/datapath/rtl/exe\_stage/}
\end{center}

The top-level execution module is:

\begin{center}
\rtl{rtl/datapath/rtl/exe\_stage/rtl/exe\_stage.sv}
\end{center}

This module acts as the dispatcher for all execution units and interfaces with the Load/Store
subsystem, scoreboards, writeback logic and pipeline control.

\begin{figure}[htbp]
  \centering
  \includegraphics[width=0.43\textwidth]{ch09_execution_overview.png}
  \caption{Execution stage overview: RR-to-EXE dispatch to Mem, ALU, Mul/Div, SIMD, FPU, and BrU}
  \label{fig:execution-overview}
\end{figure}

\section{Arithmetic Logic Unit (ALU)}

The Arithmetic Logic Unit performs all integer arithmetic and logical operations. Instead of
implementing every operation inside one large module, Sargantana divides the ALU into several
dedicated functional blocks.

The top-level ALU is implemented in:

\begin{center}
\rtl{rtl/datapath/rtl/exe\_stage/rtl/alu/alu.sv}
\end{center}

Internally, the ALU is divided into the following modules:

\begin{longtable}{@{}p{0.45\textwidth}p{0.5\textwidth}@{}}
\toprule
\textbf{RTL File} & \textbf{Function} \\
\midrule
\endhead
\rtlx{alu\_add.sv}{rtl/datapath/rtl/exe_stage/rtl/alu/alu_add.sv} & Addition and subtraction \\
\rtlx{alu\_logic.sv}{rtl/datapath/rtl/exe_stage/rtl/alu/alu_logic.sv} & AND, OR, XOR operations \\
\rtlx{alu\_shift.sv}{rtl/datapath/rtl/exe_stage/rtl/alu/alu_shift.sv} & Logical and arithmetic shifts \\
\rtlx{alu\_cmp.sv}{rtl/datapath/rtl/exe_stage/rtl/alu/alu_cmp.sv} & Comparison instructions (SLT, SLTU, etc.) \\
\rtlx{alu\_count\_pop.sv}{rtl/datapath/rtl/exe_stage/rtl/alu/alu_count_pop.sv} & Population count instructions \\
\rtlx{zero\_counter/alu\_count\_zeros.sv}{rtl/datapath/rtl/exe_stage/rtl/alu/zero_counter/alu_count_zeros.sv} & Leading/trailing zero count operations \\
\bottomrule
\end{longtable}

This modular implementation simplifies verification while allowing new ALU extensions to be
integrated independently.

\section{Branch Unit}

Branch instructions require a different execution path than arithmetic instructions because they
determine the next program counter value.

The Branch Unit performs:
\begin{itemize}
  \item Conditional branch evaluation
  \item Branch target calculation
  \item Branch taken/not-taken decision
  \item Branch misprediction recovery support
\end{itemize}

RTL implementation:

\begin{center}
\rtl{rtl/datapath/rtl/exe\_stage/rtl/branch\_unit.sv}
\end{center}

Although branch prediction occurs during the fetch stage, the Branch Unit provides the final
verification of branch outcomes during execution. Any prediction mismatch is reported back to
the pipeline so that incorrect speculative instructions can be flushed.

\section{Multiplier Unit}

Integer multiplication operations are handled by a dedicated multiplier unit.

RTL implementation:

\begin{center}
\rtl{rtl/datapath/rtl/exe\_stage/rtl/mul\_unit.sv}
\end{center}

Separating multiplication from the ALU allows multiplication instructions to execute
independently without increasing the critical path of simple arithmetic operations.

The multiplier supports RISC-V integer multiplication instructions while integrating with the
pipeline's scoreboard and scheduling logic.

\section{Divider Unit}

Division operations typically require multiple clock cycles and therefore utilize a dedicated
hardware divider.

The Divider subsystem consists of:

\begin{rtlrefs}
  \rtlitem{rtl/datapath/rtl/exe_stage/rtl/div_unit.sv}
  \rtlitem{rtl/datapath/rtl/exe_stage/rtl/div_4bits.sv}
\end{rtlrefs}

The divider performs integer division and remainder calculations while allowing the processor to
continue issuing independent instructions. The scoreboard tracks outstanding division
operations until completion.

\section{Memory Unit}

Load and store instructions are processed by the Memory Unit, which interfaces directly with the
Load/Store Queue and the cache subsystem.

Primary RTL modules include:

\begin{rtlrefs}
  \rtlitem{rtl/datapath/rtl/exe_stage/rtl/mem_unit.sv}
  \rtlitem{rtl/datapath/rtl/exe_stage/rtl/load_store_queue.sv}
  \rtlitem{rtl/datapath/rtl/exe_stage/rtl/store_buffer.sv}
  \rtlitem{rtl/datapath/rtl/exe_stage/rtl/vagu.sv}
\end{rtlrefs}

The responsibilities of these modules include:
\begin{itemize}
  \item Effective address calculation
  \item Load request generation
  \item Store buffering
  \item Memory dependency handling
  \item Interface with the memory hierarchy
\end{itemize}

The dedicated Address Generation Unit (\rtlx{vagu.sv}{rtl/datapath/rtl/exe_stage/rtl/vagu.sv}) computes effective memory addresses
before memory accesses are issued.

\section{Execution Pipeline Flow}

Once an instruction reaches the Execution stage, it is dispatched to the appropriate functional
unit according to its opcode.

The execution flow can be summarized as shown in Figure~\ref{fig:execution-flow}.

\begin{figure}[htbp]
  \centering
  \includegraphics[width=0.6\textwidth]{ch09_execution_pipeline_flow.png}
  \caption{Execution pipeline flow: dispatch through \texttt{exe\_stage.sv} to the functional units and result generation}
  \label{fig:execution-flow}
\end{figure}

The execution stage also integrates several auxiliary modules that coordinate long-latency
operations:

\begin{rtlrefs}
  \rtlitemx{pending_mem_req_queue.sv}{rtl/datapath/rtl/exe_stage/rtl/pending_mem_req_queue.sv}
  \rtlitemx{pending_fp_ops_queue.sv}{rtl/datapath/rtl/exe_stage/rtl/pending_fp_ops_queue.sv}
  \rtlitemx{score_board_scalar.sv}{rtl/datapath/rtl/exe_stage/rtl/score_board_scalar.sv}
  \rtlitemx{score_board_simd.sv}{rtl/datapath/rtl/exe_stage/rtl/score_board_simd.sv}
\end{rtlrefs}

These modules allow Sargantana to support out-of-order execution while maintaining correct
instruction dependencies.

\begin{figure}[htbp]
  \centering
  \includegraphics[width=0.28\textwidth]{ch09_execution_detail.png}
  \caption{Detailed execution-stage datapath, including the D-cache interface and auxiliary queues}
  \label{fig:execution-detail}
\end{figure}

\begin{quote}
\textit{Note: it was refined later on, I guess, because the current pipeline is a bit more advanced!}
\end{quote}
